\documentclass[12pt, titlepage]{article}
\usepackage{hyperref}

\title{Feedback Controls Lab Weekly Report \\ \normalsize{Week 4}}
\author{Aydan Vissing, Brady Schick, Gabriel Southwell, Simon Ross}
\date{\today}

\begin{document}
\maketitle

\section*{Aydan Vissing (12 hrs spent)}

This week, I started off with our team meeting on Thursday, we discussed the final pieces of the proposal, and I was assigned the Gantt chat and the budget. I spent the rest of the time there and after classes working on it and was able to get a good start. That night I also researched GitHub and how to use it so I would be properly caught up with the team on its application usage. The next day, Friday, we had our full meeting in which Dr. Fredette gave further instructions on the chart and I corrected it later with those instructions along with making it look nicer and looking over the rest of the proposal to make sure there wasn’t anything missing that we hadn’t talked about. 

\section*{Brady Schick (8.5 hrs spent)}

This week was mainly focused on preparing the project proposal. On Thursday, we met to assign the sections that still needed to be completed to individual members of the team. On Friday, we met with Dr. Fredette to get some final thoughts on the proposal. We then met to complete and finalize the proposal later that day. Throughout the week, I spent a good amount of time individually working on my sections of the proposal.

\section*{Gabriel Southwell (9.5 hrs spent)}

This week was completely overturned for me by the team presentation on Wednesday. I spent several hours preparing for that presentation and making slides. I also stopped by Dr Kohl’s office to ask for a microcontroller and met with him about a half hour. He wanted to know everything about our plans and designs so that he could give me advice on how proceed with our project. He agreed with me that based on our needs, an ESP32 is a great choice. He also lent me one for prototyping. He advised that our team should aim to integrate as much as possible onto one PCB in our final design and even in prototyping. He thinks it would be wise to meet with him as a team in the near future to talk about our plans. On my own I started writing some networking code for the ESP to test out wireless data logging.

	As a team we spent a lot of time this week working on the project proposal and talking about what to put on the gantt chart. Now that the proposal is done, the project is going to kick off in full.

\section*{Simon Ross (11.75 hrs spent)}

As Gabe mentioned, the fact that we presented from the textbook on Wednesday made this week challenging, but we were still able to accomplish a great deal on the project, with the main product being our project proposal. The proposal took up the bulk of the week, spurring many important discussions about planning and scheduling that will hopefully grease the skids for the rest of the semester. We decided on a budget, overarching schedule, and other smaller details that were included in the proposal and notes in the GitHub repository.

This upcoming week, we will focus on selecting necessary parts, developing our mathematical model, setting up our work station, and selecting a frame design.

\end{document}